% ============================================================
%  ch16  第五讲 §5.6 从几何 Satake 到经典 Satake
%  原文行 3300–3400
% ============================================================

\subsection{从几何 Satake 到经典 Satake}\label{applications of sat}

本小节讨论几何 Satake 等价与经典 Satake 同构的关系。为简单起见，假设 $G$ 在 $k=\bF_q$ 上分裂；一般情形见 \cite{RZ}。以 $\sigma$ 记 $\bF_q$ 的几何 Frobenius 自同构。固定半 Tate 扭 $\Ql(\frac{1}{2})$，等价地固定 $q$ 在 $\Ql$ 中的一个平方根 $q^{1/2}$（$=\tr(\sigma,\Ql(\frac{1}{2}))^{-1}$）。为联系几何 Satake 与经典 Satake，还需选取同构 $\iota:\Ql\cong \bC$。可安排 $\Ql(\frac{1}{2})$ 与 $\iota$ 的选取使
\[\iota(q^{1/2})\in \bR^{<0}.\]
这样选取的好处是：在 $\iota$ 下，几何 Frobenius 在 $\Ql[i](\frac{i}{2})$ 上的迹对每个 $i\in \bZ$ 都\emph{为正}。

以 $H_G$ 记经典球 Hecke 代数，即 $G(F)$ 上 $G(\mO)$-双不变的 $\bZ$-值函数全体，配以卷积积
\[(f*g)(x)=\int_{G(F)}f(y)g(y^{-1}x)dy,\]
其中测度取使 $G(\mO)$ 的体积为 $1$ 者。对 $\mu\in \xcoch(T)^+$，以 $c_\mu$ 记 $G(\mO)t^\mu G(\mO)$ 上的特征函数。则 $\{c_\mu,\ \mu\in\xcoch(T)^+\}$ 构成 $H_G$ 的一组 $\bZ$-基。

另一方面，以 $R(\hat G)$ 记 $\hat G$ 的表示环，它是以不可约表示的特征为基的 $\bZ$-代数。对 $\mu\in \xcoch(T)^+$，以 $L_\mu=\on{H}^*(\IC_\mu)$ 记相应的不可约表示，$\chi_\mu$ 记相应的特征。

经典 Satake 同构（或者更确切地说，Langlands 的重新解释）是典范同构
\[\mS: H_G\otimes\bC\cong R(\hat G)\otimes\bC.\]
细节见 \cite{Gr}。下面解释几何途径。

回忆 Grothendieck 的层-函数字典给每个 $\mA\in \Sat_G$（更一般地，$\Gr_G$ 上每个 $\ell$-进复形）配上函数
\[f_\mA: \Gr(k)\to \Ql,\quad f_\mA(x)=\sum_i(-1)^i\tr(\sigma_x, \on{H}^i_{\bar x}(\mA)),\]
其中 $\bar x$ 记 $x$ 上的几何点，$\on{H}^i_{\bar x}$ 记 $\bar x$ 处的茎上同调，$\sigma_x\in \Gal(\bar x/x)$ 记几何 Frobenius 元。因 $\mA$ 是 $L^+G$-等变的，可把 $f_\mA$ 看作 $H_G\otimes_\bZ\Ql$ 的元素。此外由定义有

\begin{lem}\label{conv vs classical conv}
对 $\mA_1,\mA_2\in\Sat_G$，有 $f_{\mA_1\star\mA_2}=f_{\mA_1}*f_{\mA_2}$。
\end{lem}

\begin{proof}
设 $m:\Gr\tilde\times\Gr\to \Gr$ 是卷积映射。给定 $x\in \Gr(k)$，纤维 $m^{-1}(x)(k)$ 可等同于 $\{(y, y^{-1}x)\mid y\in G(F)/G(\mO)\}$，且由构造有茎的典范同构 $(\mA_1\tilde\boxtimes\mA_2)_{(y,y^{-1}x)}\cong (\mA_1)_y\otimes(A_2)_{y^{-1}x}$。于是按测度的选取，
\[\begin{split}
f_{\mA_1}*f_{\mA_2}(x)&=\int_{G(F)}f_{\mA_1}(y)f_{\mA_2}(y^{-1}x)dy=\sum_{y\in G(F)/G(\mO)}f_{\mA_1}(y)f_{\mA_2}(y^{-1}x)\\
&=\sum_{y\in m^{-1}(x)(k)} \tr(\sigma_{\bar y}, (\mA_1\tilde\boxtimes\mA_2)_{\bar y})=f_{\mA_1\star\mA_2}(x).
\end{split}\]
\end{proof}

现在令 $f_\mu=f_{\IC_\mu[-(2\rho,\mu)]}$，回忆 $\IC_\mu|_{\Gr_\mu}=\Ql[(2\rho,\mu)]$。由（仿射）Schubert 簇相交上同调复形茎上同调的纯性（参见 \cite{KL}），可写
\begin{equation}\label{stalk IC}
f_\mu=c_\mu+\sum_{\la<\mu}a_{\mu\la}(q)c_\la,
\end{equation}
其中
\[a_{\mu\la}(x)=\sum_i a_{\mu\la,i}x^i\in \bN[x]\]
是 Kazhdan--Lusztig 多项式。回忆 $a_{\mu\la,i}=\dim \on{H}^{2i-(2\rho,\mu)}_{t^\la}\IC_\mu$ 是 $\IC_\mu$ 在 $t^\la$ 处次数 $2i-(2\rho,\mu)$ 的茎上同调。

回忆幺半 Abel 范畴的 $K$-群是自然的 $\bZ$-代数（即所谓的 $K$-环）。则几何 Satake 等价诱导 $\bZ$-代数同构
\[K(\Sat_G^N)\cong K(\on{Rep}_\Ql(\hat G)),\]
其中回忆 $\Sat_G^N$ 是由标准化 IC 层张成的 $\Sat_G$ 的幺半子范畴（见注 \ref{simplify}）。以 $\sigma_q$ 记 $k=\bF_q$ 的几何 Frobenius。则由 \eqref{stalk IC} 与引理 \ref{conv vs classical conv}，层-函数字典诱导代数同构
\[f: K(\Sat_G^N)\otimes\Ql\to H_G\otimes\Ql.\]
另一方面，把 $[V]$ 送到其特征 $\chi_V$ 总给出典范同构 $\on{Ch}:K(\on{Rep}_\Ql(\hat G))\cong R(\hat G)$。合起来得到同构
\[\mS':H_G\otimes\Ql\stackrel{f^{-1}}{\cong} K(\Sat_G^N)\otimes\Ql \cong K(\on{Rep}_\Ql(\hat G))\otimes\Ql\stackrel{\on{Ch}}{\cong}R(\hat G)\otimes\Ql.\]

\begin{lem}
在所选同构 $\iota:\Ql\cong\bC$ 下，$\mS'=\mS$。特别地
\[\mS^{-1}(\chi_\mu)=(-\iota(q^{1/2}))^{-(2\rho,\mu)}f_\mu.\]
\end{lem}

证明的梗概见 \cite{RZ} 末尾。首先，函子 $\on{CT}:\Sat_G\to \Sat_T$ 自然提升为函子
\[\on{CT}^N: \Sat_G^N\to \Sat_T^N.\]
（事实上由定理 \ref{MV theory} (iii) 可知 $\on{H}_c^*(\mA)$ 是 Tate 型的。）在层-函数字典与所选的 $\iota$ 下，函子 $\on{CT}^N$ 恰好对应经典 Satake 变换（例如显式公式见 \cite[\S\ 3]{Gr}）。

\begin{rmk}
(i) 若不愿引入平方根 $q^{1/2}$，可用等价 $\Sat_G^T\cong \Rep_\Ql(\hat G^T)$ 表述 Satake 同构的一个版本。

(ii) 由于对每个 $\mA\in \Sat_G^N$ 有 $\iota f_\mA\in H_G\otimes\bZ[q^{\pm 1/2}]$，同构 $\mS'$ 限制为同构 $H_G\otimes \bZ[q^{\pm 1/2}]\simeq R(\hat G)\otimes \bZ[q^{\pm 1/2}]$。若还有 $\rho\in \xch(T)$，它进一步限制为同构 $H_G\otimes\bZ[q^{-1}]\simeq R(\hat G)\otimes\bZ[q^{-1}]$。这与经典理论一致（见 \cite[\S 3]{Gr}）。
\end{rmk}

还要指出，Kazhdan--Lusztig 多项式 $a_{\mu\la}$ 可经 $\hat G$ 作如下解释。回忆命题 \ref{canonical pinning} 给出的 $\hat G$ 上的带钉，它用
\begin{equation}\label{BK fil}
N_iV=(\ker N)^{i+1}
\end{equation}
在 $\hat G$ 的任意表示 $V$ 上定义递增滤过（Brylinski--Kostant 滤过）。对 $\la\in\xcoch(T)$，以 $V(\la)$ 记 $V$ 在 $\tilde T$ 作用下的相应权子空间。则滤过 \eqref{BK fil} 诱导
\begin{equation}\label{BKfilmu}
N_iV(\la)=V(\la)\cap N_iV.
\end{equation}

下述定理在 \cite{Lu} 与 \cite{Bry} 中证明，\cite[\S~5]{Gi1} 给出另一方法（另见 \cite[\S~5]{Z11}）。

\begin{thm}
设 $\mu\in\xcoch(T)^+$，$L_\mu=\on{H}^*(\IC_\mu)$ 是 $\hat G$ 的相应表示。则
\[a_{\mu\la}(x)=\sum_i \dim \on{gr}^N_iL_\mu(\la)x^i.\]
\end{thm}
